% ============================================================================
%  06_firmware.tex — Arquitectura software del sistema integrado.
% ============================================================================
\section{Firmware}
\label{sec:firmware}

El firmware está organizado en doce módulos C++ con responsabilidades
separadas (tabla~\ref{tab:modulos}). Once de ellos heredan, sin cambios o con
adaptaciones mínimas, las responsabilidades de Act1 y Act2; el doceavo
---\texttt{diagnostics}--- es el módulo nuevo de esta entrega. El
\texttt{.ino} actúa como orquestador: invoca a cada módulo en el orden
adecuado pero no habla directamente con el hardware, manteniendo el patrón
modular que ya documentaron las dos memorias anteriores.

\begin{table}[H]
  \centering
  \small
  \begin{tabular}{p{0.30\textwidth} p{0.62\textwidth}}
    \toprule
    \textbf{Módulo} & \textbf{Responsabilidad} \\
    \midrule
    \texttt{config.h}                & Pines, \emph{setpoints}, umbrales y constantes del autodiagnóstico. \\
    \texttt{sensors.\{h,cpp\}}       & HC-SR04, DHT22 principal y \textbf{reserva}, LDR, PIR, botones y \textbf{sensor de corriente}. \\
    \texttt{actuators.\{h,cpp\}}     & Dos servos (cabina y compuerta), LEDs de control y de clasificación, buzzer. \\
    \texttt{classifier.\{h,cpp\}}    & FSM del clasificador + puente al ascensor (altura → planta). \\
    \texttt{elevator.\{h,cpp\}}      & FSM del ascensor y cola FIFO. \\
    \texttt{control.\{h,cpp\}}       & Lazos ON-OFF con histéresis y modelo térmico simulado. \\
    \texttt{ircontrol.\{h,cpp\}}     & Decodificación del mando IR (NEC). \\
    \texttt{display.\{h,cpp\}}       & HMI sobre LCD 20×4 con cuatro líneas simultáneas. \\
    \texttt{logger.\{h,cpp\}}        & Telemetría CSV unificada (clasif.\ + ascensor + clima + diagnóstico). \\
    \texttt{\textbf{diagnostics.\{h,cpp\}}} & \textbf{Capa nueva (Act3):} validación, conmutación, promediado y alarma. \\
    \texttt{actividad\_3\_vs.ino}    & \texttt{setup()} y \texttt{loop()} no bloqueantes. \\
    \bottomrule
  \end{tabular}
  \caption{Módulos del firmware integrado.}
  \label{tab:modulos}
\end{table}

\subsection{Cadena del \texttt{loop()}}

Igual que en las actividades anteriores, el bucle principal es cooperativo y
no bloqueante: cada subsistema se autorregula con \texttt{millis()}. La
diferencia respecto a Act2 es que las lecturas crudas de los sensores
ambientales ya no se hacen en el \texttt{.ino}: ahora pasan por
\texttt{diag\_tick()}, que las valida, las promedia, decide la fuente activa
y entrega al control únicamente medidas \emph{ya validadas}.

\begin{lstlisting}[caption={Cadena del \texttt{loop()} integrado.},
                   label={lst:loop}]
void loop() {
    int8_t pIR = ircontrol_poll(); if (pIR >= 0) elevator_request(pIR);
    int8_t pBT = buttonsPoll();    if (pBT >= 0) elevator_request(pBT);
    classifier_tick();   // FSM + puente al ascensor si caja ALTA
    elevator_tick();
    diag_tick();         // lee, valida, conmuta a reserva, promedia, alarma
    if (diag_hasClimate())
        control_tick(diag_temp(), diag_hum(), diag_lux());
    displayUpdate();  logger_tick();
}
\end{lstlisting}

\subsection{Integración Act1 ↔ Act2}

El módulo \texttt{classifier} encapsula la FSM heredada de Act1. Al confirmar
una caja con altura inferior a \texttt{UMBRAL\_ALTA\_CM} ejecuta el
\textbf{puente al ascensor}: mapea la altura al rango
[0, NUM_PLANTAS−1] con \texttt{map()} (caja más alta
→ planta más alta) y llama a \texttt{elevator\_request(planta)}.
Así el carril ASCENSOR de Act1 deja de ser un terminador estático y se
convierte en la entrada real al ascensor de Act2.
